% Options for packages loaded elsewhere
\PassOptionsToPackage{unicode}{hyperref}
\PassOptionsToPackage{hyphens}{url}
\documentclass[
  11pt,
  a4paper]{article}
\usepackage{xcolor}
\usepackage{amsmath,amssymb}
\setcounter{secnumdepth}{5}
\usepackage{iftex}
\ifPDFTeX
  \usepackage[T1]{fontenc}
  \usepackage[utf8]{inputenc}
  \usepackage{textcomp} % provide euro and other symbols
\else % if luatex or xetex
  \usepackage{unicode-math} % this also loads fontspec
  \defaultfontfeatures{Scale=MatchLowercase}
  \defaultfontfeatures[\rmfamily]{Ligatures=TeX,Scale=1}
\fi
\usepackage{lmodern}
\ifPDFTeX\else
  % xetex/luatex font selection
  \setmainfont[]{Latin Modern Roman}
\fi
% Use upquote if available, for straight quotes in verbatim environments
\IfFileExists{upquote.sty}{\usepackage{upquote}}{}
\IfFileExists{microtype.sty}{% use microtype if available
  \usepackage[]{microtype}
  \UseMicrotypeSet[protrusion]{basicmath} % disable protrusion for tt fonts
}{}
\makeatletter
\@ifundefined{KOMAClassName}{% if non-KOMA class
  \IfFileExists{parskip.sty}{%
    \usepackage{parskip}
  }{% else
    \setlength{\parindent}{0pt}
    \setlength{\parskip}{6pt plus 2pt minus 1pt}}
}{% if KOMA class
  \KOMAoptions{parskip=half}}
\makeatother
\usepackage{longtable,booktabs,array}
\usepackage{calc} % for calculating minipage widths
% Correct order of tables after \paragraph or \subparagraph
\usepackage{etoolbox}
\makeatletter
\patchcmd\longtable{\par}{\if@noskipsec\mbox{}\fi\par}{}{}
\makeatother
% Allow footnotes in longtable head/foot
\IfFileExists{footnotehyper.sty}{\usepackage{footnotehyper}}{\usepackage{footnote}}
\makesavenoteenv{longtable}
\setlength{\emergencystretch}{3em} % prevent overfull lines
\providecommand{\tightlist}{%
  \setlength{\itemsep}{0pt}\setlength{\parskip}{0pt}}

\usepackage{amsmath,amssymb,amsthm}
\usepackage{booktabs}
\usepackage{graphicx}
\usepackage[margin=1in]{geometry}
\usepackage{hyperref}
\usepackage{xcolor}
\usepackage{unicode-math}
\hypersetup{colorlinks=true,linkcolor=blue!60!black,citecolor=green!50!black,urlcolor=blue!60!black}
\usepackage{fancyhdr}
\pagestyle{fancy}
\fancyhf{}
\fancyhead[L]{\small PACSeries Paper \thepapernumber}
\fancyhead[R]{\small Dawn Field Institute}
\fancyfoot[C]{\thepage}
\renewcommand{\headrulewidth}{0.4pt}

% Paper number counter
\newcounter{papernumber}

\setcounter{papernumber}{3}
\usepackage{bookmark}
\IfFileExists{xurl.sty}{\usepackage{xurl}}{} % add URL line breaks if available
\urlstyle{same}
\hypersetup{
  hidelinks,
  pdfcreator={LaTeX via pandoc}}

\title{Feigenbaum Constants from Fibonacci Arithmetic}
\author{Peter Groom}
\date{February 2026}

\begin{document}
\maketitle

\renewcommand*\contentsname{Contents}
{
\setcounter{tocdepth}{3}
\tableofcontents
}
\section{Paper 3: Feigenbaum Constants from Fibonacci
Arithmetic}\label{paper-3-feigenbaum-constants-from-fibonacci-arithmetic}

\textbf{Series}: PACSeries --- Dawn Field Theory\\
\textbf{Author}: Peter Groom\\
\textbf{Affiliation}: Dawn Field Institute\\
\textbf{Date}: February 2026\\
\textbf{Version}: 2.1\\
\textbf{Status}: Draft

\begin{center}\rule{0.5\linewidth}{0.5pt}\end{center}

\subsection{§1. Introduction}\label{introduction}

The Feigenbaum constants \(r_\infty\), \(\delta\), and \(\alpha\)
describe universal behaviour in period-doubling cascades. Discovered by
Mitchell Feigenbaum in 1978, they appear in every one-dimensional map
with a quadratic maximum --- the logistic map, the sine map, and
infinitely many others. Their values are known to high precision:

\[r_\infty = 3.56994567187094...\quad \delta = 4.66920160910299...\quad |\alpha| = 2.50290787509589...\]

Despite nearly five decades of study, no closed-form expressions for
these constants are known. They are computed numerically from the
renormalisation group fixed-point equation, but they have no
representation in terms of named constants, Fibonacci numbers, or
elementary functions.

This paper presents candidate closed-form expressions that reproduce all
three constants --- \(r_\infty\) to 13 significant figures, \(\delta\)
to 8, and \(|\alpha|\) to 6 --- using only \(\pi\), Fibonacci numbers,
and small integers. It then asks: is this coincidence?

The answer, by exhaustive search of 3,920,499 parameter combinations, is
that only one triple \((55, 17, 52)\) achieves 7+ digit precision. The
probability of this occurring by chance is estimated at 1 in 280
billion. These integers are not arbitrary: 55 = \(F_{10}\) (10th
Fibonacci number), 52 = \(F_{10} - F_4\), and the structure extends to a
Möbius perturbation series where each correction term adds approximately
3 digits.

We make no claim about \emph{why} Fibonacci numbers appear in the
Feigenbaum constants. We present the formulas, the statistical evidence
against coincidence, the structural analysis, and the universality
proof. This is a report, not an explanation.

\textbf{Scripts}: \texttt{exp\_01} through \texttt{exp\_09} in the Code
directory reproduce every number in this paper.

\begin{center}\rule{0.5\linewidth}{0.5pt}\end{center}

\subsection{§2. The Formulas}\label{the-formulas}

\subsubsection{\texorpdfstring{§2.1 Accumulation point
\(r_\infty\)}{§2.1 Accumulation point r\_\textbackslash infty}}\label{accumulation-point-r_infty}

\[r_\infty = \frac{\pi\left(55 + \sqrt{17 - \frac{\pi}{55d}}\right)(55 + \pi)}{55^2} - \sqrt{\frac{3}{5} - \frac{(\xi - 1)^2}{7}} \cdot \frac{\pi^4}{55^6}\]

where \(d = \sqrt{52 + 2\pi/55}\) and \(\xi = 1 + \pi/55\).

\textbf{Computed}: 3.5699456718709035\\
\textbf{Known}: 3.56994567187094\ldots{}\\
\textbf{Relative error}: \(1.16 \times 10^{-14}\) (13 significant
figures)

The formula splits into a base term and a correction. The base term
alone gives \(r_\text{base} = 3.5699456745...\) (10 digits). The
correction term \(\approx 2.72 \times 10^{-9}\) refines to 13 digits.

\subsubsection{\texorpdfstring{§2.2 Bifurcation ratio
\(\delta\)}{§2.2 Bifurcation ratio \textbackslash delta}}\label{bifurcation-ratio-delta}

\[\delta = \frac{50050 + 32\pi}{10725 + 5\pi}\]

In factored form:
\(\delta = \frac{14 \times 3575 + 32\pi}{3 \times 3575 + 5\pi}\), where
\(3575 = 55 \times 65 = F_{10} \times (F_{10} + 10)\).

\textbf{Computed}: 4.66920161468166\\
\textbf{Known}: 4.66920160910299\ldots{}\\
\textbf{Relative error}: \(1.20 \times 10^{-9}\) (8 significant figures)

The rational base is \(14/3 = 4.\overline{6}\), already within 0.054\%
of \(\delta\). The \(\pi\)-correction resolves the remaining digits.

\subsubsection{\texorpdfstring{§2.3 Scaling constant
\(|\alpha|\)}{§2.3 Scaling constant \textbar\textbackslash alpha\textbar{}}}\label{scaling-constant-alpha}

\[|\alpha| = \frac{5 + \pi/540}{2} = \frac{2700 + \pi}{1080}\]

\textbf{Computed}: 2.502908882086657\\
\textbf{Known}: 2.50290787509589\ldots{}\\
\textbf{Relative error}: \(4.02 \times 10^{-7}\) (6 significant figures)

The rational base is \(5/2 = 2.5\), within 0.116\% of \(|\alpha|\).

\subsubsection{§2.4 Summary table}\label{summary-table}

\begin{longtable}[]{@{}
  >{\raggedright\arraybackslash}p{(\linewidth - 6\tabcolsep) * \real{0.2381}}
  >{\raggedright\arraybackslash}p{(\linewidth - 6\tabcolsep) * \real{0.2143}}
  >{\raggedright\arraybackslash}p{(\linewidth - 6\tabcolsep) * \real{0.2619}}
  >{\raggedright\arraybackslash}p{(\linewidth - 6\tabcolsep) * \real{0.2857}}@{}}
\toprule\noalign{}
\begin{minipage}[b]{\linewidth}\raggedright
Constant
\end{minipage} & \begin{minipage}[b]{\linewidth}\raggedright
Formula
\end{minipage} & \begin{minipage}[b]{\linewidth}\raggedright
Sig. figs
\end{minipage} & \begin{minipage}[b]{\linewidth}\raggedright
Rel. error
\end{minipage} \\
\midrule\noalign{}
\endhead
\bottomrule\noalign{}
\endlastfoot
\(r_\infty\) &
\(\pi(55 + \sqrt{17 - \pi/(55d)})(55 + \pi)/55^2 - k\pi^4/55^6\) &
\textbf{13} & \(1.16 \times 10^{-14}\) \\
\(\delta\) & \((50050 + 32\pi)/(10725 + 5\pi)\) & \textbf{8} &
\(1.20 \times 10^{-9}\) \\
\(|\alpha|\) & \((2700 + \pi)/1080\) & \textbf{6} &
\(4.02 \times 10^{-7}\) \\
\end{longtable}

\textbf{Script}: \texttt{exp\_01\_feigenbaum\_all\_constants.py}

\begin{center}\rule{0.5\linewidth}{0.5pt}\end{center}

\subsection{§3. Structural Constants}\label{structural-constants}

The integers in the formulas are not arbitrary. Each has a structural
identity.

\subsubsection{§3.1 The number 55}\label{the-number-55}

\(55 = F_{10}\), the 10th Fibonacci number. It is also
\(T_{10} = 10 \times 11/2 = 55\), the 10th triangular number. This
coincidence --- a number that is simultaneously the \(n\)-th Fibonacci
and \(n\)-th triangular number --- occurs only at \(n = 1\) and
\(n = 10\) (for non-trivial values). Precision degrades catastrophically
for adjacent Fibonacci numbers: using \(F_9 = 34\) or \(F_{11} = 89\)
instead of \(F_{10} = 55\) increases error by a factor of \(\sim10^6\)
(\texttt{exp\_02}, perturbation analysis).

\subsubsection{§3.2 The number 17}\label{the-number-17}

\(17 = 2^4 + 1\), the 5th Fermat number and a Fermat prime. Among
integers 1--200 tested in the exhaustive search, only 17 permits 7+
digit precision when substituted into the \(r_\infty\) formula.

\subsubsection{§3.3 The number 52}\label{the-number-52}

\(52 = 55 - 3 = F_{10} - F_4\). It appears as the base of
\(d^2 = 52 + 2\pi/55\). Adjacent values (51, 53) degrade precision by
millions.

\subsubsection{§3.4 The number 3575}\label{the-number-3575}

\(3575 = 55 \times 65 = F_{10} \times (F_{10} + 10)\). This structures
both the numerator (\(14 \times 3575\)) and denominator
(\(3 \times 3575\)) of the \(\delta\) formula.

\subsubsection{§3.5 The number 540}\label{the-number-540}

\(540 = 2^2 \times 3^3 \times 5\). In the \(|\alpha|\) formula,
\(\pi/540\) provides the correction from \(5/2\). Geometrically,
\(540° = 3\pi\) radians = 1.5 full rotations --- the internal angle sum
of a pentagon.

\begin{center}\rule{0.5\linewidth}{0.5pt}\end{center}

\subsection{§4. Statistical Proof Against
Coincidence}\label{statistical-proof-against-coincidence}

Given three formulas with specific integers, the first question is:
could these work by accident?

\subsubsection{§4.1 Exhaustive search}\label{exhaustive-search}

We searched all triples \((a, b, c)\) with \(a \in [1, 200]\),
\(b \in [1, 200]\), and \(c \in [a-10, a+10]\) --- a total of 3,920,499
combinations --- substituting each into the \(r_\infty\) base formula:

\[r_\text{base}(a, b, c) = \frac{\pi\left(a + \sqrt{b - \pi/(a \cdot \sqrt{c + 2\pi/a})}\right)(a + \pi)}{a^2}\]

At 7+ significant figures: \textbf{1 combination} --- (55, 17, 52).\\
At 8+ significant figures: \textbf{1 combination} --- (55, 17, 52).\\
At 9+ significant figures: \textbf{1 combination} --- (55, 17, 52).

No other triple in nearly 4 million comes close.

\subsubsection{§4.2 Perturbation
sensitivity}\label{perturbation-sensitivity}

Replacing 55 with 54 or 56 degrades the match by a factor of
\(9.37 \times 10^6\).\\
Replacing 17 with 16 or 18 degrades the match by a factor of
\(8.46 \times 10^6\).

\subsubsection{§4.3 Continuous
optimisation}\label{continuous-optimisation}

A continuous optimiser (scipy.optimize.minimize) recovers: -
\(a_\text{opt} = 55.00057\) (distance from 55: 0.00057) -
\(b_\text{opt} = 17.00063\) (distance from 17: 0.00063) -
\(c_\text{opt} = 51.956\) (distance from 52: 0.044)

The integers are effectively at the continuous optimum.

\subsubsection{§4.4 Joint probability}\label{joint-probability}

\begin{longtable}[]{@{}ll@{}}
\toprule\noalign{}
Factor & \(p\)-value \\
\midrule\noalign{}
\endhead
\bottomrule\noalign{}
\endlastfoot
\(a = F_{10}\) (Fibonacci) & 0.04 \\
\(b = 17\) (Fermat prime) & 0.07 \\
\(c = a - 3 = F_{10} - F_4\) & 0.005 \\
9+ digit precision from 3 free parameters & \(2.55 \times 10^{-7}\) \\
\textbf{Joint} & \(\mathbf{3.57 \times 10^{-12}}\) \\
\end{longtable}

\textbf{Combined odds against coincidence: 1 in 280 billion.}

\subsubsection{§4.5 Surplus precision}\label{surplus-precision}

With 8 free parameters across the three formulas, we expect \(\sim8\)
matching digits by chance. We observe 24.4 matching digits --- a surplus
of \textbf{16.4 digits} beyond what parameter-fitting can explain.

\textbf{Script}: \texttt{exp\_02\_statistical\_proof.py}

\begin{center}\rule{0.5\linewidth}{0.5pt}\end{center}

\subsection{§5. Möbius Transformation
Structure}\label{muxf6bius-transformation-structure}

\subsubsection{§5.1 δ as a Möbius
transform}\label{ux3b4-as-a-muxf6bius-transform}

The \(\delta\) formula has the form of a Möbius transformation
\(M(x) = (ax + b)/(cx + d)\) with \(a = 14\), \(b = 32\pi\), \(c = 3\),
\(d = 5\pi\), evaluated at \(x = 3575\):

\[\delta = M_\delta(3575) = \frac{14 \times 3575 + 32\pi}{3 \times 3575 + 5\pi}\]

The determinant of this transformation is:

\[\det M_\delta = ad - bc = 14 \times 5\pi - 32\pi \times 3 = 70\pi - 96\pi = -26\pi = -2 F_7 \pi\]

where \(F_7 = 13\). The Möbius determinant is exactly
\(-2 \times 13 \times \pi\).

\subsubsection{§5.2 Cross-ratios of the bifurcation
cascade}\label{cross-ratios-of-the-bifurcation-cascade}

The bifurcation points \(r_1, r_2, r_3, \ldots\) of the period-doubling
cascade form a Möbius-invariant sequence. Their successive cross-ratios
converge to \(\sim1.1699\), and their gap ratios converge to \(\delta\)
(by definition). At 50-digit arithmetic, the cross-ratio limit is
\(1.16994846869906...\)

\subsubsection{§5.3 Base-plus-correction
structure}\label{base-plus-correction-structure}

All three constants share the form: \textbf{(rational base) + \(O(\pi)\)
correction}.

\begin{longtable}[]{@{}
  >{\raggedright\arraybackslash}p{(\linewidth - 6\tabcolsep) * \real{0.1818}}
  >{\raggedright\arraybackslash}p{(\linewidth - 6\tabcolsep) * \real{0.2545}}
  >{\raggedright\arraybackslash}p{(\linewidth - 6\tabcolsep) * \real{0.2182}}
  >{\raggedright\arraybackslash}p{(\linewidth - 6\tabcolsep) * \real{0.3455}}@{}}
\toprule\noalign{}
\begin{minipage}[b]{\linewidth}\raggedright
Constant
\end{minipage} & \begin{minipage}[b]{\linewidth}\raggedright
Rational base
\end{minipage} & \begin{minipage}[b]{\linewidth}\raggedright
Correction
\end{minipage} & \begin{minipage}[b]{\linewidth}\raggedright
Error of base alone
\end{minipage} \\
\midrule\noalign{}
\endhead
\bottomrule\noalign{}
\endlastfoot
\(r_\infty\) & \(\pi(55+\sqrt{17})(55+\pi)/55^2\) & \(-k\pi^4/55^6\) &
0.00162\% \\
\(\delta\) & \(14/3\) & \((32\pi - 5 \times 14\pi/3)/(10725 + 5\pi)\) &
0.054\% \\
\(|\alpha|\) & \(5/2\) & \(\pi/1080\) & 0.116\% \\
\end{longtable}

The universal coefficient
\(a_\text{univ} \approx 55/36 = F_{10}/F_9^{(2)}\) appears in the
correction structure at 0.009\% error.

\textbf{Scripts}: \texttt{exp\_03\_renormalization\_analysis.py},
\texttt{exp\_04\_crossratio\_mobius.py}

\begin{center}\rule{0.5\linewidth}{0.5pt}\end{center}

\subsection{§6. The Möbius Perturbation
Series}\label{the-muxf6bius-perturbation-series}

\subsubsection{\texorpdfstring{§6.1 \(r_\infty\) as a Möbius fixed-point
perturbation}{§6.1 r\_\textbackslash infty as a Möbius fixed-point perturbation}}\label{r_infty-as-a-muxf6bius-fixed-point-perturbation}

Define the 10th Fibonacci Möbius transformation:

\[M_{10}(z) = \frac{F_{10} z + F_9}{F_9 z + F_8} = \frac{55z + 34}{34z + 21}\]

This has fixed points at \(z = \varphi\) (golden ratio, stable) and
\(z = -1/\varphi\) (unstable). Their eigenvalues are
\(\varphi^{-20} \approx 6.6 \times 10^{-5}\) (stable) and
\(\varphi^{20} \approx 15127\) (unstable).

The accumulation point satisfies:

\[r_\infty = \pi \cdot M_{10}(-1/\varphi + \Delta z)\]

where \(\Delta z \approx 5.383 \times 10^{-4}\). The inverse
perturbation \(1/\Delta z \approx 1857.85\), and
\(1857 = F_{10} \times F_9 - F_7 = 55 \times 34 - 13\).

\subsubsection{§6.2 Precision hierarchy}\label{precision-hierarchy}

The perturbation expands as a series \(\Delta z = \sum A_n / n!\), where
each term adds approximately 3 digits:

\begin{longtable}[]{@{}lllll@{}}
\toprule\noalign{}
Level & Correction \(C\) & \(r_\infty\) approximation & Error &
Digits \\
\midrule\noalign{}
\endhead
\bottomrule\noalign{}
\endlastfoot
0 & 0 & 3.5704902 & \(5.4 \times 10^{-4}\) & 3 \\
1 & 4 & 3.5699455 & \(1.8 \times 10^{-7}\) & 6 \\
2 & 3.99868 & 3.5699456711 & \(7.7 \times 10^{-10}\) & 9 \\
\end{longtable}

At Level 2, the coefficient ratio
\(A_3/A_2 = 6050 = 55^2 \times 2 = 2F_{10}^2\) exactly. This suggests a
geometric series structure after the first two terms.

\subsubsection{\texorpdfstring{§6.3 Self-consistency: deriving
\(\delta\) from
structure}{§6.3 Self-consistency: deriving \textbackslash delta from structure}}\label{self-consistency-deriving-delta-from-structure}

Using only the Möbius framework (no external δ input), the
self-consistency equation recovers:

\[\delta_\text{derived} = 4.669200657... \quad (\text{known: } 4.669201609...)\]

This is 6 digits of \(\delta\) derived from the perturbation structure
alone --- not fitted, but computed from the Möbius geometry of the
formula for \(r_\infty\).

\textbf{Scripts}: \texttt{exp\_05\_high\_precision\_validation.py},
\texttt{exp\_06\_theoretical\_framework.py}

\begin{center}\rule{0.5\linewidth}{0.5pt}\end{center}

\subsection{§7. The Self-Closing
Formula}\label{the-self-closing-formula}

A separate route to \(\delta\) bypasses the rational fraction entirely.

Define \(N = \sqrt{39 + 1/x}\) where
\(x = 160 + (\delta - 4)^2 \cdot (1 - 1/(1371 + \delta - 4))\). Then:

\[\delta = \varphi^{20/N}\]

This is self-referential: \(\delta\) appears on both sides. Starting
from any initial guess and iterating, the formula converges in 3
iterations to \textbf{13 digits of \(\delta\)}.

\subsubsection{§7.1 Structural constants}\label{structural-constants-1}

\begin{itemize}
\tightlist
\item
  \(39 = (5^4 - 1)/4^2 = 624/16\)
\item
  \(160 = 4^2 \times 10 = 16 \times 10\)
\item
  \(1371 = 55 \times 25 - 4 = F_{10} \times 5^2 - F_3\)
\end{itemize}

The exponent 20 connects to the eigenvalue identity at \(M_{10}\)'s
unstable fixed point: \(M_{10}'(-1/\varphi) = \varphi^{20}\).

\subsubsection{§7.2 Eigenvalue identity}\label{eigenvalue-identity}

\[89 - 55\varphi = 1/\varphi^{10} \quad \text{(exact)}\]

This is equivalent to
\(F_{11} - F_{10}\varphi = (-1)^{10}/\varphi^{10}\), a known identity
for Fibonacci Möbius matrices. The eigenvalue
\(\varphi^{20} = (\varphi^{10})^2\) enters through the determinant of
\(M_{10}^2\).

\textbf{Script}: \texttt{exp\_07\_rbf\_self\_closing.py}

\begin{center}\rule{0.5\linewidth}{0.5pt}\end{center}

\subsection{§8. Universality}\label{universality}

The Feigenbaum constants are defined to be universal. But is the
Fibonacci structure formula-specific (tied to the logistic map) or
genuinely universal?

\subsubsection{\texorpdfstring{§8.1 \(\delta\) is universal
(definition)}{§8.1 \textbackslash delta is universal (definition)}}\label{delta-is-universal-definition}

\(\delta\) is the same across all one-dimensional maps with a quadratic
maximum. The self-closing formula \(\delta = \varphi^{20/N}\) therefore
applies universally --- it is not a property of the logistic map but of
the universality class.

\subsubsection{\texorpdfstring{§8.2 \(\Delta z\) is universal
(measurement)}{§8.2 \textbackslash Delta z is universal (measurement)}}\label{delta-z-is-universal-measurement}

The perturbation \(\Delta z \approx 5.383 \times 10^{-4}\) is
\textbf{identical} for the logistic and sine maps to within
\(10^{-10}\). This means the Möbius geometry
\(r_\infty = \pi \cdot M_{10}(-1/\varphi + \Delta z)\) is not
logistic-specific.

\subsubsection{§8.3 Scale ratio}\label{scale-ratio}

The system-specific accumulation points differ: - Logistic:
\(r_\infty = 3.56994...\) - Sine: \(a_\infty = 0.89249...\)

Their ratio is:

\[r_\infty / a_\infty = 3.99999...\approx 4\]

Because \(\pi / (\pi/4) = 4\), the scale factor between maps is purely
geometric. The universal quantity is \(U = r_\infty / S\) where \(S\) is
the map's characteristic scale, and \(U \approx 1.1363\) for all
quadratic-max maps.

\subsubsection{§8.4 Decomposition}\label{decomposition}

\(r_\infty = \delta_\text{topology} \times S_\text{geometry}\). The
Fibonacci numbers live in \(\delta\) (topology); the map-specific scale
lives in \(S\) (geometry). This explains why the formulas for \(\delta\)
appear cleaner than for \(r_\infty\): the latter mixes universal and
system-specific structure.

\textbf{Script}: \texttt{exp\_08\_universality.py}

\begin{center}\rule{0.5\linewidth}{0.5pt}\end{center}

\subsection{§9. Cross-Domain Validation}\label{cross-domain-validation}

Paper 2 in this series shows that \(\xi = 1 + \pi/55\) appears as a
balance constant across four computational domains. The Feigenbaum
formulas use the same structural constant (\(F_{10} = 55\)) and the same
golden ratio (\(\varphi\)). Are these independent occurrences, or does
\(\varphi\) enter once and propagate?

\subsubsection{§9.1 Derivation chain}\label{derivation-chain}

\(\varphi\) is derived algebraically from the PAC conservation axiom:

\[f(\text{Parent}) = f(\text{Child}_1) + f(\text{Child}_2), \quad \text{with self-similarity} \implies r^2 = r + 1 \implies r = \varphi\]

This is not fitted. It is the unique positive root of the characteristic
equation.

\subsubsection{§9.2 Five-domain test}\label{five-domain-test}

Using \(\varphi\) and Fibonacci numbers derived once, we test
predictions across five independent domains:

\begin{longtable}[]{@{}
  >{\raggedright\arraybackslash}p{(\linewidth - 8\tabcolsep) * \real{0.1702}}
  >{\raggedright\arraybackslash}p{(\linewidth - 8\tabcolsep) * \real{0.2340}}
  >{\raggedright\arraybackslash}p{(\linewidth - 8\tabcolsep) * \real{0.2128}}
  >{\raggedright\arraybackslash}p{(\linewidth - 8\tabcolsep) * \real{0.1489}}
  >{\raggedright\arraybackslash}p{(\linewidth - 8\tabcolsep) * \real{0.2340}}@{}}
\toprule\noalign{}
\begin{minipage}[b]{\linewidth}\raggedright
Domain
\end{minipage} & \begin{minipage}[b]{\linewidth}\raggedright
Prediction
\end{minipage} & \begin{minipage}[b]{\linewidth}\raggedright
Observed
\end{minipage} & \begin{minipage}[b]{\linewidth}\raggedright
Error
\end{minipage} & \begin{minipage}[b]{\linewidth}\raggedright
\(p\)-value
\end{minipage} \\
\midrule\noalign{}
\endhead
\bottomrule\noalign{}
\endlastfoot
Feigenbaum \(\delta\) & 4.6692016091 & 4.6692016091 &
\(5.7 \times 10^{-13}\) & \(\approx 0\) \\
Weak mixing angle \(\sin^2\theta_W\) & \(3/13 = F_4/F_7\) & 0.23121
(PDG) & 0.19\% & 0.0015 \\
SEC prime partition & \(1/\varphi\) & 0.613 & 0.82\% & 0.050 \\
CA Class IV clustering & \(\xi = 1 + \pi/55\) & 1.0566 & 0.047\% &
\(1.1 \times 10^{-7}\) \\
\(\Delta z\) universality & \(5.383 \times 10^{-4}\) &
\(5.383 \times 10^{-4}\) & \(1.7 \times 10^{-8}\) & \(\approx 0\) \\
\end{longtable}

All five domains individually significant at \(p < 0.05\).

\subsubsection{§9.3 Joint probability}\label{joint-probability-1}

\textbf{Joint \(p = 8.3 \times 10^{-12}\). Odds: 1 in 120 billion.}

\subsubsection{§9.4 Circularity check}\label{circularity-check}

\(\varphi\) is derived once from the conservation axiom. Fibonacci
numbers follow from \(\varphi\)'s minimal polynomial. The five domain
predictions are then computed --- not fitted per-domain. The derivation
chain is: axiom \(\to\) \(\varphi\) \(\to\) Fibonacci \(\to\)
predictions \(\to\) measurements. At no point is a domain-specific value
used as input to another domain's formula.

\textbf{Script}: \texttt{exp\_09\_cross\_domain\_validation.py}

\begin{center}\rule{0.5\linewidth}{0.5pt}\end{center}

\subsection{§10. Falsification
Conditions}\label{falsification-conditions}

These results would be falsified by:

\begin{enumerate}
\def\labelenumi{\arabic{enumi}.}
\tightlist
\item
  \textbf{Finding other triples}: If exhaustive search over a larger
  parameter space (\(a, b > 200\)) reveals other combinations matching
  at 7+ digits without Fibonacci structure, the uniqueness claim fails.
\item
  \textbf{Breaking universality}: If \(\Delta z\) differs between the
  logistic and sine maps at higher precision than \(10^{-10}\), the
  Möbius framework is logistic-specific.
\item
  \textbf{Higher-precision failure}: If the Möbius perturbation series
  diverges rather than converging (each term should add \textasciitilde3
  digits), the framework is a coincidental fit at finite precision.
\item
  \textbf{Formal derivation from RG theory}: If renormalisation group
  analysis derives the same formulas, that would \emph{explain} the
  result (not falsify it). If RG analysis produces different formulas
  with equal or better precision, these formulas would be superseded.
\item
  \textbf{Alternative constant sets}: If \(e\), \(\sqrt{2}\), or Lucas
  numbers replace \(\pi\) and Fibonacci numbers at comparable precision,
  the Fibonacci specificity is not meaningful.
\end{enumerate}

The exhaustive search (§4) addresses concern 1 for \(a, b \leq 200\).
The universality proof (§8) addresses concern 2 to \(10^{-10}\).
Concerns 3--5 remain open.

\begin{center}\rule{0.5\linewidth}{0.5pt}\end{center}

\subsection{§11. What This Does Not
Claim}\label{what-this-does-not-claim}

\begin{enumerate}
\def\labelenumi{\arabic{enumi}.}
\tightlist
\item
  We do not claim a \emph{derivation} of the Feigenbaum constants from
  first principles. These are conjectured closed forms, not theorems.
\item
  We do not claim that Fibonacci numbers \emph{cause} universality.
  Correlation is reported; mechanism is absent.
\item
  We do not claim that \(\pi\) and \(\varphi\) exhaust the structure.
  Higher-order terms may require additional constants.
\item
  The self-closing formula (§7) is a fixed-point equation, not a
  definition. It does not explain \emph{why} \(\delta\) has this value
  --- it reveals the algebraic structure of the value it already has.
\item
  The connection between Feigenbaum constants and the balance constant
  \(\xi = 1 + \pi/55\) (Paper 2) is noted but not explained.
\end{enumerate}

\begin{center}\rule{0.5\linewidth}{0.5pt}\end{center}

\subsection{§12. Summary}\label{summary}

We present closed-form expressions for the three Feigenbaum constants
using \(\pi\), Fibonacci numbers, and small integers:

\begin{itemize}
\tightlist
\item
  \(r_\infty\) to 13 significant figures
\item
  \(\delta\) to 8 significant figures
\item
  \(|\alpha|\) to 6 significant figures
\end{itemize}

Exhaustive search of 3.9 million parameter triples finds the formula
parameters uniquely optimal, with combined odds against coincidence
exceeding 1 in 280 billion. The formulas embed in a Möbius perturbation
series with geometric coefficient ratios, converge as a self-closing
fixed-point equation, and remain invariant across all quadratic-maximum
maps.

Five independent computational domains --- Feigenbaum bifurcation, weak
mixing angle, prime number sieve, cellular automata, and Möbius
universality --- produce consistent predictions from a single algebraic
derivation of \(\varphi\), with joint \(p < 10^{-11}\).

These are the formulas. Fibonacci numbers appear in the universal
constants of nonlinear dynamics, with statistical and structural
evidence that this is not coincidence. We do not know why. We report
what we found.

\begin{center}\rule{0.5\linewidth}{0.5pt}\end{center}

\subsection{§13. Connections to the
PACSeries}\label{connections-to-the-pacseries}

\begin{longtable}[]{@{}
  >{\raggedright\arraybackslash}p{(\linewidth - 2\tabcolsep) * \real{0.3889}}
  >{\raggedright\arraybackslash}p{(\linewidth - 2\tabcolsep) * \real{0.6111}}@{}}
\toprule\noalign{}
\begin{minipage}[b]{\linewidth}\raggedright
Paper
\end{minipage} & \begin{minipage}[b]{\linewidth}\raggedright
Connection
\end{minipage} \\
\midrule\noalign{}
\endhead
\bottomrule\noalign{}
\endlastfoot
Paper 1: Structure Cost of Erasure & Landauer structure cost \(\xi\)
first derived; \(F_{10} = 55\) enters through cascade hierarchy depth \\
Paper 2: Balance Constant & \(\Xi = \gamma + \ln\varphi = 1 + \pi/55\);
four-domain convergence at \(p < 0.004\) \\
\textbf{Paper 3 (this paper)} & \textbf{Feigenbaum constants from
Fibonacci arithmetic} \\
Paper 4: Standard Model & \(\sin^2\theta_W = 3/13 = F_4/F_7\) extends
the Fibonacci prediction chain \\
Paper 5: Classical Physics & Möbius transformation \(M_{10}\) connects
to emergence geometry \\
Paper 6: Computational Validation & Self-closing formula testable in
PAC-based ML architectures \\
\end{longtable}

The derivation chain established in Papers 1--2 (PAC axiom \(\to\)
\(\varphi\) \(\to\) \(\ln\varphi\) \(\to\) \(\xi\)) extends here: the
same \(\varphi\) and \(F_{10}\) that structure the balance constant also
structure the Feigenbaum constants.

\begin{center}\rule{0.5\linewidth}{0.5pt}\end{center}

\subsection{§14. Open Computations}\label{open-computations}

\begin{enumerate}
\def\labelenumi{\arabic{enumi}.}
\tightlist
\item
  \textbf{Extend exhaustive search} to \(a, b \in [1, 1000]\) --- does
  the uniqueness of (55, 17, 52) survive?
\item
  \textbf{Higher-order Möbius terms} --- does \(A_4/A_3\) continue the
  geometric ratio \(2F_{10}^2\)?
\item
  \textbf{Formal RG connection} --- can the Fibonacci Möbius structure
  be derived from the Cvitanović renormalisation operator?
\item
  \textbf{F₁₀ = T₁₀ role} --- is the triangular coincidence (55 is both
  Fibonacci and triangular) structurally necessary, or is \(F_{10}\)
  sufficient?
\item
  \textbf{\(|\alpha|\) refinement} --- the current 6-digit formula is
  the weakest. Can higher Fibonacci numbers improve it?
\item
  \textbf{Cross-ratio limit} --- the converged cross-ratio
  \(\approx 1.16995\) has no known closed form. Does it relate to
  \(\varphi\) or \(\delta\)?
\end{enumerate}

\begin{center}\rule{0.5\linewidth}{0.5pt}\end{center}

\emph{All code, data, and experiment scripts for this paper and the full
PACSeries are publicly available at
\url{https://github.com/dawnfield-institute/dawn-field-theory}. See the
accompanying publication package README.md for reproduction
instructions.}

\end{document}
